% sections/07_conclusion.tex
\section{Conclusion}
\label{sec:conclusion}

This report compared four supervised models on the \texttt{penguins} dataset, two for classification and two for regression, and measured how much the sex of an individual adds to predictions based on physical traits alone. Both studies reached high performance on the test set. Logistic regression classified species with an accuracy of up to 0.985 and a weighted precision of 0.986, and the decision tree reached 0.955 on both metrics. For regression, the random forest predicted body mass with an R\textsuperscript{2} of up to 0.841, while linear regression reached 0.786. Hyperparameter optimization of the decision tree did not improve on the default configuration, which indicates that the tree was already well suited to this dataset under the metric used. The main finding is that adding sex to the predictors produces a consistent gain in every model where the effect is measurable: classification with logistic regression gained between 0.045 and 0.060 in accuracy, and regression gained between 0.048 and 0.075 in R\textsuperscript{2}, always with the same sign across two disjoint training subsets.

Three limitations should be considered when reading the results. First, the decision tree was left unconstrained in the final comparison. Although cross-validation selected the default parameters, an unlimited depth is a structural risk on datasets with more noise or less separation between classes, and a future iteration should test the same pipeline on such data. Second, the MinMaxScaler was fitted on the full cleaned frame before the train/test split. This is a methodological leakage: the scaler saw the test rows during fitting, which in this case has a small numerical effect but violates the rule that no information from the test set should influence training. The correct fix is to wrap the scaler and the model in a \texttt{Pipeline} from \texttt{scikit-learn} so that the scaler is refitted inside each fold. Third, the test set has only 67 rows. The standard error of an accuracy estimate at that size is close to 0.03, so the magnitudes of the reported gains should be treated as approximate. The sign of the effect of sex is stable across training subsets, which supports the qualitative conclusion, but the exact numbers may shift with a different random split.

The results of this study are reproducible. The pipeline uses a fixed random state, the A/B split is deterministic, and the notebooks and the report sources are versioned in a public repository. The approach can be extended to other datasets with the same structure, and the A/B design in particular is a low-cost way to check whether an observed effect is real or an artifact of the training sample.